% Plantilla CV Actuarial — ActuaryLab
% Enfoque: distintiva pero ATS-friendly (una columna, sin tablas complejas ni foto)
% Uso sugerido: Overleaf con compilador pdfLaTeX
% Sin instalación local necesaria — sube este archivo y compila directamente
\documentclass[10pt,letterpaper]{article}
\usepackage[margin=1.35cm]{geometry}
\usepackage[utf8]{inputenc}
\usepackage[T1]{fontenc}
\usepackage[spanish]{babel}
\usepackage[hidelinks]{hyperref}
\usepackage{enumitem}
\usepackage{titlesec}
\usepackage{xcolor}
\usepackage{lmodern}

\pagenumbering{gobble}
\setlength{\parindent}{0pt}
\setlist[itemize]{leftmargin=*, itemsep=2pt, topsep=2pt}

\definecolor{navy}{HTML}{0D1F2D}
\definecolor{sage}{HTML}{6BA08D}
\definecolor{softgray}{HTML}{666666}

\titleformat{\section}
  {\large\bfseries\color{navy}}
  {}{0em}{}[\vspace{-2pt}{\color{sage}\titlerule[0.8pt]}]
\titlespacing*{\section}{0pt}{9pt}{5pt}

\newcommand{\cvitem}[2]{\textbf{#1:} #2\\}
\newcommand{\role}[4]{
  \textbf{#1} \hfill {\color{softgray} #2}\\
  \textit{#3}\\[-2pt]
  #4\vspace{3pt}
}

\begin{document}

% Sin foto: recomendado para versiones digitales y ATS.
{\Huge \textbf{Nombre Apellido}}\\[-1pt]
{\color{sage}\rule{\textwidth}{1.2pt}}\\[3pt]
\href{mailto:nombre.apellido@email.com}{nombre.apellido@email.com} \; | \; Toluca, México \; | \; \href{https://linkedin.com/in/tu-nombre}{linkedin.com/in/tu-nombre} \; | \; \href{https://github.com/tuusuario}{github.com/tuusuario}

\section*{Perfil}
Estudiante de Actuaría / Actuario(a) junior con interés en ciencias actuariales, seguros, regulación y análisis de datos. Experiencia aplicada en R/Python, documentación reproducible y proyectos con datos públicos.

\section*{Experiencia y proyectos aplicados}
\role{Servicio Social · Facultad de Economía, UAEMéx}{2025--2026}{R/Python · análisis demográfico · artículo científico}{
\begin{itemize}
  \item Procesé, limpié y validé bases de datos demográficos para su aplicación en un artículo científico utilizando R/Python.
  \item Elaboré resúmenes estadísticos y visualizaciones para apoyar la interpretación de resultados.
\end{itemize}}

\role{Proyecto actuarial · Graduación de tablas de mortalidad}{2026}{R · CUSF Anexo 6.3.8 · documentación reproducible}{
\begin{itemize}
  \item Realicé la graduación de tablas de mortalidad con datos públicos aplicando criterios regulatorios del Anexo 6.3.8 de la CUSF.
  \item Documenté supuestos, metodología y resultados en un reporte reproducible.
\end{itemize}}

\role{Club de Actuaría UAEMéx · Área de habilidades blandas}{2024--2025}{Comunicación profesional · organización estudiantil}{
\begin{itemize}
  \item Coordiné actividades de integración y comunicación profesional para estudiantes de actuaría.
  \item Apoyé en difusión, logística y organización de talleres y eventos académicos.
\end{itemize}}

\section*{Educación}
\textbf{Licenciatura en Actuaría} \hfill 2022--2026\\
Universidad / Facultad

\section*{Habilidades técnicas}
\cvitem{R}{tidyverse, ggplot2, ajuste de distribuciones, simulación Monte Carlo, reportes reproducibles}
\cvitem{Python}{pandas, numpy, análisis exploratorio, visualización, scikit-learn básico}
\cvitem{Datos y documentación}{SQL básico, Git/GitHub, Quarto, README técnico}
\cvitem{Regulación}{LISF, CUSF, CNSF, EPA; comprensión inicial de LCNBV si el perfil se orienta a banca o inversiones}
\cvitem{Cloud / BI}{Power BI/Tableau básico; AWS/Azure nivel inicial si aplica al puesto}

\section*{Formación complementaria}
\begin{itemize}
  \item Cursos, congresos, seminarios o comunidades actuariales relevantes.
  \item Para perfiles de banca, inversiones o riesgos financieros: considerar formación complementaria como AMIB Figura 3 si aplica al objetivo profesional.
\end{itemize}

\end{document}
